% Imported from LaTeX_and_pdfs/AGAT_dump.tex; see provenance.json.
\part{Laplace Transform}


\begin{quotation}
This is for you, you, my number one \\
This is for you, you, my number one \\

This is for you, you, my number one \\

Spotlights, big stage \\
Fifty thousand fans screamin' in a rage \\
Bodyguards and limousines \\
This is the way I see you in my dreams \\
\end{quotation}
- Usher, "Superstar" \\

\begin{quotation}
Me: In which category or functor in category theory do Laplace transforms belong in? \\

Chat-GPT: Laplace transforms can be studied in the category of topological vector spaces. More specifically, the Laplace transform can be seen as a functor from the category of functions of exponential growth to the category of holomorphic functions defined on a half-plane in the complex plane. \\

In this category, the morphisms are continuous linear maps that preserve the exponential growth of functions. The Laplace transform functor takes a function in the source category and maps it to its Laplace transform in the target category. \\

The use of category theory provides a powerful framework for studying Laplace transforms and their properties in a more abstract and general way.
\end{quotation}

\section{Laplace Transform: Definitions and Properties}

\subsection{Laplace Transform Definition and Examples}

We will follow Mauch (2004) \cite{agat:Mauc2004}.

\begin{definition}
\begin{equation}
\mathcal{L}[f(t)] = \int_0^{\infty} e^{-st} f(t) dt \equiv \widehat{f}(s)
\end{equation}
\end{definition}

\begin{definition}
A function $f(t)$ is of \textbf{exponential order} $\alpha$ if $\exists \, $ constants $t_0$, and $M$ s.t.
\begin{equation}
|f(t) | < M e^{\alpha t} \quad \, \forall \, t > t_0
\end{equation}
\end{definition}

If $\exists \, $ $\int_0^{t_0} f(t) dt$ and $f(t)$ is of exponential order $\alpha$, then $\exists \, \widehat{f}(s)$ \, $\forall \, \Re{(s)} > \alpha$.

\begin{proof}
\[
\begin{gathered}
| \int_0^{\infty} e^{-st} f(t) dt | = | \int_0^{\infty} \exp{ (-\Re{(s)} t - i \Im{ (s)} t )} f(t) dt | \leq \int_0^{\infty} \exp{( -\Re{(s)} t)} |f(t)| dt < \\
< \int_0^{\infty} \exp{ (- \Re{(s)} t) } Me^{\alpha t} dt   = M \int_0^{\infty} dt \exp{ (- (\Re{ (s)} - \alpha ) t ) } = M \left. \left( \frac{ \exp{ (- (\Re{ (s)} - \alpha ) t} }{ - (\Re{ (s)} - \alpha ) } \right) \right|_0^{\infty} = \frac{M }{ \Re{ (s) } - \alpha }
\end{gathered}
\]
if $\Re{(s)} > \alpha$
\end{proof}

For example, \\
$\sin{(t)}$ is of exponential order 0. \\
$te^{2t}$ is of exponential order $\alpha$, \, $\forall \, \alpha >2$.

Example 31.1 (Mauch (2004) \cite{agat:Mauc2004}). Consider the Laplace transform of $f(t) = 1$.  Since $f(t)= 1$ is of exponential order $\alpha $ \, $\forall \, \alpha >0$, Laplace transform integral converges for $\Re{(s)} < 0$.

\[
\begin{gathered}
\widehat{f}(s) = \int_0^{\infty} e^{-st} dt = \left[ \frac{-1}{s} e^{-st} \right]_0^{\infty} = \frac{1}{s}
\end{gathered}
\]

Example 31.1.3. Consider Laplace transform of the Heaviside function
\[
\begin{gathered}
H(t-c) = \begin{cases}
0 & \text{ for } t < c \\
1 & \text{ for } t > c
\end{cases}
\end{gathered}
\]
where $c>0$.

\[
\begin{gathered}
\mathcal{L}[H(t-c)] = \int_0^{\infty} e^{-st} H(t-c) dt = \int_c^{\infty} e^{-st} dt = \left[ \frac{e^{-st}}{-s} \right]_c^{\infty} = \frac{e^{-cs}}{s} \quad \, \forall \, \Re{(s)} > 0
\end{gathered}
\]

Example 31.1.4. Consider $H(t-c) f(t-c)$
\[
\begin{gathered}
\mathcal{L}[H(t-c) f(t-c)] = \int_0^{\infty} dt e^{-st} H(t-c) f(t-c) = \int_c^{\infty} dt e^{-st} f(t-c) = \int_0^{\infty} dt e^{-s(t+c) } f(t) = e^{-cs} \widehat{f}(s)
\end{gathered}
\]

\subsubsection{Laplace Transform Properties and Formulae}

\begin{theorem}
\begin{equation}
\begin{aligned}
& \mathcal{L}\left[ \frac{d}{dt} f(t) \right] = s\widehat{f}(s) - f(0) \\
& \mathcal{L}\left[ \frac{d^2}{dt^2} f(t) \right] = s^2 \widehat{f}(s) - sf(0) - f'(0)
\end{aligned}
\end{equation}
\end{theorem}

\begin{proof}
To deserve these formulas,
\[
\begin{gathered}
\begin{aligned}
\mathcal{L}\left[ \frac{d}{dt} f(t) \right] & = \int_0^{\infty} dt e^{-st} f'(t) = \left[ e^{-st}f(t) \right]_0^{\infty} - \int_0^{\infty}dt (-se^{-st}) f(t) = -f(0) + s\widehat{f}(s) \\
\mathcal{L}\left[ \frac{d^2}{dt^2} f(t) \right] & = \int_0^{\infty} dt e^{-st} f''(t) = \left[ e^{-st} f'(t) \right]_0^{\infty} - \int_0^{\infty} dt (-se^{-st} ) f'(t) = -f'(0) + s \mathcal{L} [ \frac{d}{dt} f'(t) ] = \\
& = -f'(0) - sf(0 ) + s^2 \widehat{f}(s)
\end{aligned}
\end{gathered}
\]
\end{proof}

\subsection{Inverse Laplace Transform: Definition}

\begin{definition}
The inverse Laplace transform is denoted
\begin{equation}
f(t) = \mathcal{L}^{-1}[ \widehat{f}(s) ]
\end{equation}
Compute the inverse Laplace transform with the Mellin inversion formula:

\begin{equation}
f(t) = \frac{1}{2\pi i} \int_{\alpha - i \infty}^{\alpha + i \infty} ds e^{st} \widehat{f}(s)
\end{equation}
where $\alpha$ is a real constant that's to the right of the singularities of $\widehat{f}(s)$.
\end{definition}

Consider 
\[
\begin{gathered}
\mathcal{L} [ \mathcal{L}^{-1} [\widehat{f}(s)] ] =  \mathcal{L} \left[ \frac{1}{2\pi i} \int_{\alpha - i \infty}^{\alpha + i \infty} dz e^{zt} \widehat{f}(z) \right] = \int_0^{\infty} dt e^{-st} \frac{1}{2\pi i} \int_{\alpha - i \infty}^{\alpha + i \infty} dz e^{zt} \widehat{f}(z)
\end{gathered}
\]
Interchange order of integration:
\[
\begin{gathered}
\frac{1}{2\pi i} \int_{\alpha - i \infty }^{\alpha + i \infty} dz \widehat{f}(z) \int_0^{\infty} dt e^{-(s- z) t}
\end{gathered}
\]
Since $\Re{ (z)} = \alpha$ (because $z$ is from $\alpha - i \infty$ to $\alpha + i \infty$, the integral in $t$ exists for $\Re{ (s)} > \alpha$.
So
\[
\begin{gathered}
\int_0^{\infty} dt e^{-(s-z)t} =  \left( \frac{ e^{-(s-z)t} }{ -(s-z) } \right)_0^{\infty} = \frac{1}{s-z} \Longrightarrow \frac{1}{2 \pi i} \int_{\alpha - i \infty}^{\alpha + i \infty} dz \frac{\widehat{f}(z) }{ s-z}
\end{gathered}
\]

